\section{实践指南：构建高效 Agent Harness 的原则}

基于 LangChain 的实验和 deepagents 的整体构建经验，以下是构建 Agent Harness 的核心原则。

\subsection{原则一：代替 Agent 做上下文组装}

当前 Agent 在陌生环境中的上下文组装能力还不成熟。主动注入以下信息可以大幅减少可避免的错误：

\begin{itemize}[leftmargin=1.5em]
    \item 目录结构和可用工具
    \item 编码最佳实践
    \item 问题解决策略框架
    \item 评估标准和约束条件
\end{itemize}

\subsection{原则二：激进地推动 Self-Verification}

模型对自己的第一个方案有天然的偏好。必须通过 Prompt 和 Middleware 双管齐下，推动 Agent 主动验证：
\begin{itemize}[leftmargin=1.5em]
    \item 运行测试
    \item 对照任务规格（而非自身代码）检查结果
    \item 迭代改进直到通过验证
\end{itemize}

\subsection{原则三：用 Trace 作为反馈信号}

Trace 是 Agent 自我评估和调试的基础。关键是同时调试\textbf{工具}和\textbf{推理}——模型走错路的原因往往是缺少某个工具或不知道如何使用它。

\subsection{原则四：短期内检测并修复坏模式}

模型目前还不完美。Harness 设计者的职责是：
\begin{itemize}[leftmargin=1.5em]
    \item 围绕当前模型的短板做设计（如循环检测、强制验证）
    \item 同时为更强模型规划演进路径
    \item 接受这些护栏最终会被淘汰——但今天它们是必要的
\end{itemize}

\subsection{原则五：为每个模型定制 Harness}

通用 Harness 是一个好的起点，但最优性能需要针对具体模型做迭代：
\begin{itemize}[leftmargin=1.5em]
    \item 跑几轮 Trace 分析 $\rightarrow$ 改进循环
    \item 关注模型特有的失败模式
    \item 利用模型特有的优势
\end{itemize}

\subsection{本章小结}

五条原则构成了一个完整的 Harness Engineering 方法论：上下文注入降低环境不确定性、Self-Verification 确保输出质量、Trace 分析驱动迭代改进、护栏设计应对当前模型短板、模型专属定制最大化性能。


